% ============================================================
%  MACROS PERSONNALISEES
% ============================================================

% --- Indentation de bloc pour sections ---
\newenvironment{sectionblock}{%
  \begin{adjustwidth}{1.5cm}{0cm}%
}{%
  \end{adjustwidth}%
}

% --- Valeurs et incertitudes (siunitx-friendly) ---
\newcommand{\val}[2]{\SI{#1}{#2}}
\newcommand{\valinc}[3]{\SI{#1 +- #2}{#3}}
\newcommand{\valrel}[4]{\SI{#1 +- #2}{#3}\ (\SI{#4}{\percent})}
\newcommand{\urel}[2]{\SI{#1}{#2}}
\newcommand{\ecartRel}[1]{\SI{#1}{\percent}}
\newcommand{\zscore}[1]{\ensuremath{z=\num{#1}}}

% --- Equations avec legende ---
\newcommand{\eqleg}[2]{%
  \begin{equation}\label{#1}
    #2
  \end{equation}
  \noindent\small\textit{Legende (Eq.~\ref{#1}) :}\quad
}

% --- References intelligentes ---
\newcommand{\figref}[1]{\cref{fig:#1}}
\newcommand{\tabref}[1]{\cref{tab:#1}}
\newcommand{\secref}[1]{\cref{sec:#1}}

% --- Boite equation inline (style sobre) ---
\newtcbox{\eqinline}{
  on line,
  enhanced,
  colback=black!4,
  colframe=black!30,
  sharp corners,
  boxrule=0.5pt,
  left=6pt,
  right=6pt,
  top=2pt,
  bottom=2pt,
  nobeforeafter,
  tcbox raise base
}
\newcommand{\eqbox}[1]{\eqinline{$\displaystyle #1$}}

% --- Boites d'information ---
\newtcolorbox{boxinfo}{
  colback=BoxBg,
  colframe=BoxBorder,
  boxrule=0.85pt,
  left=7pt,
  right=7pt,
  top=6pt,
  bottom=6pt,
  arc=0.6mm
}

% --- Marqueurs de travail ---
\newcommand{\todo}[1]{\textcolor{HighlightRed}{\textbf{[TODO : #1]}}}
\newcommand{\exemple}[1]{{\color{HighlightBlueVivid}\itshape #1}}
\newcommand{\fictif}[1]{\textcolor{HighlightBlueVivid}{#1}}

% --- Resultats encadres ---
\newenvironment{resultkey}[2]{%
  \begin{tcolorbox}[
    colback=HeigBlueLight,
    colframe=BoxBorder,
    boxrule=0.9pt,
    arc=0.6mm,
    title=\textbf{#1},
    fonttitle=\sffamily\bfseries,
    coltitle=white,
    colbacktitle=BoxBorder,
    breakable
  ]
    \large $\boxed{#2}$
    \par\vspace{4pt}
}{%
  \end{tcolorbox}
}

% --- Environnements pour info et attention ---
\newtcolorbox{info}{
  colback=BoxBg,
  colframe=BoxBorder,
  boxrule=0.85pt,
  left=7pt,
  right=7pt,
  top=6pt,
  bottom=6pt,
  title=Remarque,
  fonttitle=\sffamily\bfseries,
  coltitle=white,
  colbacktitle=BoxBorder,
  arc=0.6mm,
  breakable
}

\newtcolorbox{attention}{
  colback=HighlightOrangeVivid!8,
  colframe=HighlightOrangeVivid,
  boxrule=0.9pt,
  left=7pt,
  right=7pt,
  top=6pt,
  bottom=6pt,
  title=Attention,
  fonttitle=\sffamily\bfseries,
  coltitle=white,
  colbacktitle=HighlightOrangeVivid,
  arc=0.6mm,
  breakable
}

% --- Compatibilite anciennes macros ---
\newcommand{\valeur}[2]{\SI{#1}{#2}}

% --- Macros orientees rapport d'ingenierie ---
\newcommand{\req}[2]{\noindent\textbf{Exigence #1.}\ #2\par}
\newcommand{\assump}[2]{\noindent\textbf{Hypothese #1.}\ #2\par}
\newcommand{\decision}[2]{\noindent\textbf{Decision #1.}\ #2\par}
\newcommand{\riskitem}[2]{\noindent\textbf{Risque #1.}\ #2\par}

\newcommand{\resultatcle}[3]{%
  \begin{tcolorbox}[colback=HeigBlueLight, colframe=BoxBorder, boxrule=0.9pt, arc=0.6mm]
    \textbf{Resultat cle :} #1\\
    \textbf{Valeur :} #2\\
    \textbf{Confiance/limite :} #3
  \end{tcolorbox}
}

\newcommand{\assetref}[2]{\texttt{#1} (\textit{#2})}
\newcommand{\tracecmd}[1]{\texttt{#1}}
